% Options for packages loaded elsewhere
\PassOptionsToPackage{unicode}{hyperref}
\PassOptionsToPackage{hyphens}{url}
\documentclass[
  ignorenonframetext,
]{beamer}
\newif\ifbibliography
\usepackage{pgfpages}
\setbeamertemplate{caption}[numbered]
\setbeamertemplate{caption label separator}{: }
\setbeamercolor{caption name}{fg=normal text.fg}
\beamertemplatenavigationsymbolsempty
% remove section numbering
\setbeamertemplate{part page}{
  \centering
  \begin{beamercolorbox}[sep=16pt,center]{part title}
    \usebeamerfont{part title}\insertpart\par
  \end{beamercolorbox}
}
\setbeamertemplate{section page}{
  \centering
  \begin{beamercolorbox}[sep=12pt,center]{section title}
    \usebeamerfont{section title}\insertsection\par
  \end{beamercolorbox}
}
\setbeamertemplate{subsection page}{
  \centering
  \begin{beamercolorbox}[sep=8pt,center]{subsection title}
    \usebeamerfont{subsection title}\insertsubsection\par
  \end{beamercolorbox}
}
% Prevent slide breaks in the middle of a paragraph
\widowpenalties 1 10000
\raggedbottom
\AtBeginPart{
  \frame{\partpage}
}
\AtBeginSection{
  \ifbibliography
  \else
    \frame{\sectionpage}
  \fi
}
\AtBeginSubsection{
  \frame{\subsectionpage}
}
\usepackage{iftex}
\ifPDFTeX
  \usepackage[T1]{fontenc}
  \usepackage[utf8]{inputenc}
  \usepackage{textcomp} % provide euro and other symbols
\else % if luatex or xetex
  \usepackage{unicode-math} % this also loads fontspec
  \defaultfontfeatures{Scale=MatchLowercase}
  \defaultfontfeatures[\rmfamily]{Ligatures=TeX,Scale=1}
\fi
\usepackage{lmodern}
\ifPDFTeX\else
  % xetex/luatex font selection
\fi
% Use upquote if available, for straight quotes in verbatim environments
\IfFileExists{upquote.sty}{\usepackage{upquote}}{}
\IfFileExists{microtype.sty}{% use microtype if available
  \usepackage[]{microtype}
  \UseMicrotypeSet[protrusion]{basicmath} % disable protrusion for tt fonts
}{}
\makeatletter
\@ifundefined{KOMAClassName}{% if non-KOMA class
  \IfFileExists{parskip.sty}{%
    \usepackage{parskip}
  }{% else
    \setlength{\parindent}{0pt}
    \setlength{\parskip}{6pt plus 2pt minus 1pt}}
}{% if KOMA class
  \KOMAoptions{parskip=half}}
\makeatother
\setlength{\emergencystretch}{3em} % prevent overfull lines
\providecommand{\tightlist}{%
  \setlength{\itemsep}{0pt}\setlength{\parskip}{0pt}}
\usepackage{bookmark}
\IfFileExists{xurl.sty}{\usepackage{xurl}}{} % add URL line breaks if available
\urlstyle{same}
\hypersetup{
  hidelinks,
  pdfcreator={LaTeX via pandoc}}

\author{\texorpdfstring{}{}}
\date{}

\begin{document}

\begin{frame}[fragile]{Methods}
\protect\phantomsection\label{methods}
The \emph{Docxology Template} architecture is bifurcated into a globally
shared \texttt{infrastructure/} layer and project-specific
\texttt{projects/} silos. This ``Two-Layer Architecture'' allows for
massive scaling across heterogeneous research domains (e.g., from
cognitive diagrams to code optimization) while maintaining a singular
source of truth for build logic.

\begin{block}{The 7-Stage Pipeline}
\protect\phantomsection\label{the-7-stage-pipeline}
Our primary methodology involves a sequential, staged orchestrator that
manages the transition between code and prose:

\begin{enumerate}
\tightlist
\item
  \textbf{Stage 00 (Sanitization)}: Environment validation and
  dependency syncing.
\item
  \textbf{Stage 01 (Verification)}: Full execution of
  \texttt{tests/infra\_tests} and
  \texttt{projects/\textless{}name\textgreater{}/tests}.
\item
  \textbf{Stage 02 (Extraction)}: Triggering analysis scripts in
  \texttt{projects/\textless{}name\textgreater{}/scripts/} to generate
  \texttt{data/} and \texttt{figures/}.
\item
  \textbf{Stage 03 (Rendering)}: Sequential Pandoc and XeLaTeX
  invocation to compile Markdown sections into a unified PDF.
\item
  \textbf{Stage 04 (Security)}: Application of SHA-256/512 hashing and
  Alpha-Channel steganographic text/QR overlays.
\item
  \textbf{Stage 05 (Validation)}: Structural PDF checking for
  xref/trailer integrity.
\item
  \textbf{Stage 06 (Summarization)}: Generation of Executive Reports and
  LLM-aided reviews.
\end{enumerate}
\end{block}

\begin{block}{High-Consistency Documentation}
\protect\phantomsection\label{high-consistency-documentation}
We utilize a ``Triad Documentation'' standard for all infrastructure
components:

\begin{itemize}
\tightlist
\item
  \texttt{README.md}: High-level purpose and quick-start.
\item
  \texttt{AGENTS.md}: Deep technical implementation details for AI
  collaborators.
\item
  \texttt{SKILL.md}: Actionable patterns and API usage snippets.
\end{itemize}
\end{block}
\end{frame}

\end{document}
